\chapter{Conclusion}
\noindent
This thesis set out to answer the question:

\begin{quote}
How can a DevSecOps toolchain be designed to support ETL workflows in regulated
industries, enhancing auditability and data governance?
\end{quote}

The answer developed in this thesis is that DevSecOps can support governed ETL
when the delivery pipeline and the data pipeline are treated as one evidence
system. In the proposed architecture, a dataset is not only the output of an
Airflow run. It is the result of a governed change, an approved release, a
specific infrastructure state, a known code bundle, a quality decision, and a
recorded lineage path. Auditability is therefore designed into the workflow
rather than reconstructed after the workflow has already completed.

\section{Summary of Contributions}
\noindent
The first contribution of this thesis is the architecture itself. Chapter~5
presented a governed ETL architecture at escalating levels of detail, beginning
with the high-level system structure and moving through process maps, data
flows, metadata flows, and a worked execution path. This architecture extends a
standard Airflow and AWS-style ETL pipeline with change governance, release
manifests, policy checks, quality gates, lineage publication, metadata
records, and audit-chain artifacts.

The second contribution is the implemented prototype. The implementation used a
governed EDGAR pipeline to show how the architecture can be materialised in
code. The prototype produced run manifests, release references, change
references, code hashes, quality results, OpenLineage-style events,
OpenMetadata-compatible records, and audit-chain files. Although the
demonstration was intentionally scoped, it showed that a produced dataset
version can be traced back to the authority, code, infrastructure, and controls
that produced it.

The third contribution is the evaluation framework and requirement mapping.
Chapter~3 translated auditability and governance concerns into five system
requirements: governance accountability, information security assurance,
operational risk control, continuous control monitoring, and end-to-end
evidence chain. Chapter~7 then evaluated the prototype against these
requirements and linked each one back to the auditability dimensions of
normativity, accountability, referentiality, reliability, and verifiability.

\section{Findings}
\noindent
The main finding is that auditability improves when evidence is captured as
part of the delivery and runtime workflow. The baseline pipeline could produce
data and operational logs, but it did not bind the output to a complete
governance chain. The proposed architecture changed this by making the run
manifest and metadata artifacts central to the system.

The worked example in Chapter~6 demonstrated this point. Starting from the
example dataset version, the evidence chain could be followed to the
audit-chain artifact, run manifest, release manifest, Terraform reference,
change record, code hashes, quality decision, lineage event, and metadata
records. This does not remove the need for human judgement in an audit, but it
does change what the auditor is judging. Instead of assembling evidence from
unrelated logs and repository history, the auditor can inspect a structured
chain of fields and artifacts.

The evaluation also showed that the architecture is strongest when controls
produce evidence, not only decisions. A policy check that blocks deployment is
useful, but it becomes auditable when its result is retained and linked to the
release or run. A data contract is useful, but it becomes auditable when the
quality result records that the contract was checked before promotion. A
lineage record is useful, but it becomes stronger when it points to the same
dataset version and manifest as the rest of the evidence chain.

\section{Scope and Reflection}
\noindent
The demonstration was intentionally scoped around the evidence model. The EDGAR
pipeline and representative run were used to test whether a governed ETL system
could connect release identity, runtime state, quality evidence, lineage, and
metadata into a coherent audit chain. Production-scale reliability, multi-team
access-control behaviour, and long-term retention are operational extensions of
that same architecture rather than the focus of the evaluation.

The most important design trade-off is reliance on GitHub. GitHub is useful
because it brings source control, issues, pull requests, CI/CD, and release
automation close together. At the same time, that convenience concentrates
trust. The outages encountered during development made this risk clear. A
future production version of this architecture should split some of these
responsibilities across independent services so that audit evidence is not
dependent on one platform remaining continuously available and internally
consistent.

\section{Closing Remarks}
\noindent
This thesis demonstrates that DevSecOps can enhance ETL governance when it is
used to create durable evidence, not only automation. The proposed architecture
does not treat compliance as a final review step. It treats governance
references, release identity, code integrity, quality decisions, lineage, and
metadata as normal outputs of the pipeline.

For regulated environments, this matters because trust in data depends on more
than correct transformation logic. It depends on whether the organisation can
explain how the data was produced, who authorised the change, which controls
ran, and how the output can be reconstructed. The architecture developed in
this thesis provides a practical path toward that goal by making auditability a
designed property of the ETL system.
